% 2011 年山东卷理科第 9 题：四个选项的函数图象（选项图）。
% 来源 PDF：普通高考/2011/2011山东理.pdf，SHA-256 d0f409fd3fa82b12c210d587c8412d4c446946a201f87d124e2fc17c5a50db16。
% 原页：第 1 页（A3 三栏）第 1 栏下部接第 2 栏上部，书页码 689。
% 重绘依据：原页四图为示意图，纵轴「4」刻度与横轴「2π」刻度的实际比例与 y=x/2-2sin x
%   的取值并不一致（原页纵向压缩），故按 600 dpi 逐像素提取曲线拓扑后重绘，刻度位置照原页：
%   (A) 曲线形如 y=0.31x-1.24cos x（含偶函数项，不关于原点对称，左支与右支不同号）；
%   (B) 曲线形如 y=0.61x-1.84sin x（奇函数，振幅相对一次项小，|x|>4 出图）；
%   (C) 曲线形如 y=0.31x-1.26sin x，与 y=(x/2-2sin x) 同形（零点 ±2.475，极值 x≈±1.32、±4.97、±7.60）；
%   (D) 为奇函数曲线，小极值在 x≈±1.5、大极值在 x≈±6.6，按原页数字化控制点平滑重绘。
%   四图的坐标轴、箭头、刻度「4」「2π」与「O」的位置均照原页，不按比例。
\begin{tikzpicture}[>=stealth, semithick, font=\small, x=0.17cm, y=0.28cm]

  % ================= (A) =================
  \begin{scope}
    \draw[->] (-12.1,0) -- (12.2,0) node[right] {$x$};
    \draw[->] (0,-4.35) -- (0,4.30) node[above] {$y$};
    \draw (-2.4,4) -- (2.4,4);
    \node[left] at (-2.6,4) {$4$};
    \draw (6.283,-0.22) -- (6.283,0.22);
    \node[below] at (6.283,-0.22) {$2\pi$};
    \node at (-2.4,-1.5) {$O$};
    \draw[thick, domain=-10.1:10.1, samples=280]
      plot (\x, {0.31*\x - 1.24*cos(\x r)});
    \node at (-6.4,-6.2) {(A)};
  \end{scope}

  % ================= (B) =================
  \begin{scope}[xshift=7.6cm]
    \draw[->] (-12.1,0) -- (12.2,0) node[right] {$x$};
    \draw[->] (0,-4.35) -- (0,4.30) node[above] {$y$};
    \draw (-2.4,4) -- (2.4,4);
    \node[left] at (-2.6,4) {$4$};
    \draw (6.283,-0.22) -- (6.283,0.22);
    \node[below] at (6.283,-0.22) {$2\pi$};
    \node at (-2.4,-1.5) {$O$};
    \draw[thick, domain=-4.1:4.1, samples=200]
      plot (\x, {0.61*\x - 1.84*sin(\x r)});
    \node at (-6.4,-6.2) {(B)};
  \end{scope}

  % ================= (C) =================
  \begin{scope}[yshift=-6.6cm]
    \draw[->] (-12.1,0) -- (12.2,0) node[right] {$x$};
    \draw[->] (0,-4.35) -- (0,4.30) node[above] {$y$};
    \draw (-2.4,4) -- (2.4,4);
    \node[left] at (-2.6,4) {$4$};
    \draw (6.283,-0.22) -- (6.283,0.22);
    \node[below] at (6.283,-0.22) {$2\pi$};
    \node at (-2.4,-1.5) {$O$};
    \draw[thick, domain=-10.1:10.1, samples=280]
      plot (\x, {0.31*\x - 1.26*sin(\x r)});
    \node at (-6.4,-6.2) {(C)};
  \end{scope}

  % ================= (D) =================
  \begin{scope}[xshift=7.6cm, yshift=-6.6cm]
    \draw[->] (-12.1,0) -- (12.2,0) node[right] {$x$};
    \draw[->] (0,-4.35) -- (0,4.30) node[above] {$y$};
    \draw (-2.4,4) -- (2.4,4);
    \node[left] at (-2.6,4) {$4$};
    \draw (6.283,-0.22) -- (6.283,0.22);
    \node[below] at (6.283,-0.22) {$2\pi$};
    \node at (-2.4,-1.5) {$O$};
    \draw[thick] plot[smooth, tension=0.75] coordinates {
      (-10.08,3.42) (-9.58,2.00) (-8.57,-0.70) (-7.57,-2.38) (-6.55,-2.78)
      (-5.55,-2.20) (-4.53,-1.21) (-3.53,-0.11) (-2.52,0.90) (-1.52,1.30)
      (-0.50,0.70) (0.15,0.20) (0,0) (-0.15,-0.20) (0.50,-0.71)
      (1.52,-1.31) (2.52,-0.90) (3.53,0.11) (4.53,1.21) (5.55,2.20)
      (6.55,2.78) (7.56,2.39) (8.57,0.68) (9.58,-2.02) (10.08,-3.42)};
    \node at (-6.4,-6.2) {(D)};
  \end{scope}

\end{tikzpicture}
